\section{Symmetric Groups}

\begin{definition}
    Given a set \(X\), a permutation of \(X\) is
    a bijection \(X \rightarrow X\)\\
    \(S_X=\left\{ \sigma : X \rightarrow X \mid \sigma \t{ is a bijection } \right\}\)
\end{definition}

\begin{lemma}
    For all sets \(X \neq \emptyset\), \(S_X\) is a group
\end{lemma}

\begin{proof}
    \phantom{}

    \begin{itemize}
        \item Closure: Suppose \(\sigma\t{ , }\mu \in S_X\)\\
        Consider \(\sigma \circ \mu\) since it is a composition of bijections
        it is also a bijection.\\ Also functions go from
        \(X \rightarrow X\) thus \(\sigma \circ \mu\) is the same.\\
        Therefore, the composition is a bijection that is mapped from 
        \(X \rightarrow X\) thus \(\sigma \circ \mu \in S_X\)
        \item Identity: Consider \(e : X \rightarrow X\) s.t. \(e(x)\mapsto x\)\\
        Thus \(e \in S_x\) and for all \(\sigma \in S_X\;\; \sigma \circ e(x)=\sigma(x)\)\\ 
        Thus \(\sigma \circ e = \sigma\). Also, it follows that \(e \circ \sigma=\sigma\)
        \item Inverses: Since \(S_X\) has only bijective functions
        all bijective functions have inverses.
        \item Associativity: We have associativity by the fact that function
        composition has associativity.
    \end{itemize}
\end{proof}

\begin{definition}
    The permutation group on \(n\)-objects \(S_n\) is the \(S_X\) with \(X=[n]\)\\
    Observe \(\mid S_n \mid = n!\)
\end{definition}

\begin{example}
    \(\sigma\),\(\mu\in S_4\)
    \[\sigma=\left( 
    \begin{array}{cccc}
    1 & 2 & 3 & 4 \\
    2 & 3 & 4 & 1 \\
    \end{array}
    \right)
    \mu=\left( 
    \begin{array}{cccc}
    1 & 2 & 3 & 4 \\
    3 & 4 & 2 & 1 \\
    \end{array}
    \right)
    \]
    %
    \[\sigma^{-1}=\left( 
    \begin{array}{cccc}
    1 & 2 & 3 & 4 \\
    4 & 1 & 2 & 3 \\
    \end{array}
    \right)
    \mu^{-1}=\left( 
    \begin{array}{cccc}
    1 & 2 & 3 & 4 \\
    4 & 3 & 1 & 2 \\
    \end{array}
    \right)
    \]
    %
    \[\sigma\mu=\left( 
    \begin{array}{cccc}
    1 & 2 & 3 & 4 \\
    4 & 1 & 3 & 2 \\
    \end{array}
    \right)
    \mu\sigma=\left( 
    \begin{array}{cccc}
    1 & 2 & 3 & 4 \\
    4 & 2 & 1 & 3 \\
    \end{array}
    \right)
    \]
\end{example}
\newpage
\subsection{Cycle Notation}
\begin{example}
    \(\sigma,\mu\in S_5\)
    \[\sigma=\left( 
    \begin{array}{ccccc}
    1 & 2 & 3 & 4 & 5 \\
    3 & 4 & 2 & 5 & 1\\
    \end{array}
    \right)
    \mu=\left( 
    \begin{array}{ccccc}
    1 & 2 & 3 & 4 & 5\\
    4 & 5 & 2 & 1 & 3\\
    \end{array}
    \right)
    \]
    \[
    \sigma=\left( 13245 \right)
    \mu=\left( 14 \right)\left( 253 \right)
    \]

    Example Calculations:
    \begin{align*}
        \sigma\mu &=\left( 13245 \right)\left( 14 \right)\left( 253 \right)\\
                  &=\left( 152 \right)\left( 34 \right)\\
        \mu\sigma &=\left( 14 \right)\left( 253 \right)\left( 13245 \right)\\
                  &=\left( 12 \right)\left( 354 \right)\\
        \mu^2     &=\left( 14 \right)\left( 253 \right)\left( 14 \right)\left( 253 \right)\\
                  &=\left( 14 \right)^2\left( 253 \right)^2=\left( 235 \right)
    \end{align*}
\end{example}
\begin{lemma}
    Disjoint cycles are commuative
\end{lemma}
\begin{proof}
    Let \(\alpha,\beta\) be disjoint cycles of \([n]\)

    Let \(\alpha=\left( a_1,a_2,\dots,a_k \right)\), and \(\beta=\left( b_1,b_2,\dots,b_j \right)\)

    Let \(x\in[n]\)

    Consider: \(\alpha\beta(x)\),\(\beta\alpha(x)\)

    If \(x\in\alpha\) then \(x\notin\beta \implies \beta(x)=x\implies\alpha\beta(x)=\alpha(x)=a_i\)
                      
    \hspace*{2.28in}  \(\implies \beta\alpha(x)=\beta(a_i)=a_i\) since \(a_i \notin \beta\)

    Use similar reasoning for if \(x\in\beta\)

    \(\implies \alpha\beta=\beta\alpha\)
\end{proof}

\begin{lemma}
    Every \(\sigma\in S_n\) can be written as disjoint cycles
\end{lemma}

\begin{lemma}
    The order of an \(m\)-cylce is \(m\)
\end{lemma}
\newpage

\begin{lemma}
    Suppose \(G\) is a group and ord\((g)=n\). Then \(g^m=e\iff n\mid m\)
\end{lemma}

\begin{proof}
    Let \(G\) be a group and ord\((g)=n\)\\
    (\(\Longleftarrow\))

    Assume \(n\mid m\)

    \(\implies\) \(m=nk\), \(k\in\Z\)

    \(g^m=g^nk=(g^n)^k=(e)^k=e\)\\
    (\(\implies\))

    Assume \(g^m=e\) and use division algorithm to write \(m=nq+r\) s.t. \(0\leq r<n\)

    \(e=g^m=g^{nq+r}=(g^n)^qg^r=g^r \implies r=0 \implies n\mid m\)
\end{proof}

\begin{theorem}
    If \(\sigma_1,\sigma_2,\dots,\sigma_k\) are disjoint cycles of length
    \(a_1,a_2,...,a_k\) and\\ \(\sigma=\sigma_1,\sigma_2,\dots,\sigma_k\).
    Then ord\((\sigma)=\t{lcm}(a_1,a_2,\dots,a_k)\).
\end{theorem}

\begin{proof}
    Let \(m=\)ord\((\sigma)\), \(l=\t{lcm}(a_1,a_2,\dots,a_k)\)

    Since \(l\) is a lcm it is also a multiple so \(l=a_jq_j\) for \(1\leq j\leq k\) where \(q_j\in\Z\)
    
    \begin{align*}
        \sigma^l&=\sigma_1^l\sigma_2^l\cdots\sigma^l & \t{We can do this because of Lemma 2.1.6}\\
                &=(\sigma_1^{a_1})^{q_1}\cdots(\sigma_k^{a_k})^{q_k}\\
                &=(1) &\t{Apply Lemma 2.1.8}
    \end{align*}

    \(\implies m\mid l\)

    Also, \((1)=\sigma^m=\sigma_1^m\cdots\sigma_k^m\) since they are disjoint none of them are inverses of each other
    
    \(\implies\) \(\sigma_1^m=\cdots=\sigma_k^m=(1)\)

    \(\t{(Lemma 2.1.9)}\implies\) \(a_1\mid m \cdots a_k \mid m \implies l \mid m\)

    Thus, since \(l \mid m\) and \(m \mid l\) \(m=l\)
\end{proof}
\newpage
\begin{example}
    \phantom{text}

    \begin{itemize}
        \item Find an element of order 12 in \(S_7\)
        
        \((1234)(567)\in S_7\) by Theorem 2.0.10 \(\implies\) the order is \(3\cdot4=12=\)lcm(3,4)
        \item Possible orders of elements in \(S_5\)
        \begin{itemize}
            \item 1: (1), 5=1+1+1+1+1
            \item 2: (12), 5=2+1+1+1+1 or (12)(34), 5=2+2+1
            \item 3: (123), 5=3+1+1
            \item 4: (1234), 5=4+1
            \item 5: (12345), 5=5
            \item 6: (12)(345), 5=2+3
        \end{itemize}
    \end{itemize}
\end{example}
